% Part: incompleteness
% Chapter: theories-computability
% Section: inseparability

\documentclass[../../../include/open-logic-section]{subfiles}

\begin{document}

\olfileid{inc}{tcp}{ins}

\olsection{$\Th{Q}$માં સાબિત અને નકારી શકાય એવાં વાક્યો
  સંગણનીય રીતે અવિભાજ્ય છે}


$\Th{\bar Q}$ એ એવાં
વાક્યોનો ગણ માનો કે જેના
{\em નિષેધો} $\Th{Q}$માં
સાબિત થાય છે; એટલે
$\Th{\bar Q} = \Setabs{!A}{\Th{Q} \Proves
  \lnot !A}$.
યાદ રાખો કે વિચ્છિન્ન
ગણો $A$ અને $B$ને
સંગણનીય રીતે અવિભાજ્ય
ત્યારે કહીએ જ્યારે એવો
સંગણનીય ગણ~$C$ ન હોય
કે $A \subseteq C$ અને
$B \subseteq \Complement{C}$.

\begin{lem}
$\Th{Q}$ અને
$\Th{\bar Q}$ સંગણનીય
રીતે અવિભાજ્ય છે.
\end{lem}

\begin{proof}
માનો કે $C$ એવો સંગણનીય
ગણ છે કે $\Th{Q} \subseteq C$
અને $\Th{\bar Q} \subseteq \Complement{C}$.
$R(x,y)$ સંબંધ નીચે મુજબ લો:
\begin{quote}
$x$ એ સૂત્ર $!D(u)$નો
સંકેત છે અને $!D(\num y)$
એ $C$માં છે.
\end{quote}
$R(x,y)$ સાર્વત્રિક
સંગણનીય સંબંધ છે એમ
બતાવીને વિરોધાભાસ
મેળવીશું.

માનો કે $S(y)$ સંગણનીય
છે અને તેને $\Th{Q}$માં
$!D_S(u)$ નિરૂપે છે.
ત્યારે
\begin{eqnarray*}
S(\num n) & \lif & \Th{Q} \Proves !D_S(\num n) \\
& \lif & !D_S(\num n) \in C
\end{eqnarray*}
અને
\begin{eqnarray*}
\lnot S(\num n) & \lif & \Th{Q} \Proves \lnot !D_S(\num n) \\
& \lif & !D_S(\num n) \in \Th{\bar Q} \\
& \lif & !D_S(\num n) \not\in C
\end{eqnarray*}
આથી $S(y)$ એ
$R(\#(!D_S(u)),y)$ની
સમકક્ષ છે.
\sourcecorrection{OLINC-041}{સ્થિર અંગ્રેજી
મૂળના અંતિમ સંકેતમાં
$!D_S(\num u)$ લખાયું છે.
અહીં $R$નો પહેલો ઘટક
મુક્ત ચલવાળા સૂત્ર
$!D_S(u)$નો સંકેત છે;
ચલને અંકરૂપ આપતું
ચિહ્ન તેથી દૂર કર્યું છે.}
\end{proof}

\end{document}
